\originalchapter{8.04 作业1及独立解答}
\label{ch:ps1}
本章题面按 \href{https://ocw.mit.edu/courses/8-04-quantum-physics-i-spring-2016/resources/mit8_04s16_ps1_2016/}{8.04 Spring 2016 Problem Set 1} 中文转写, 发布于2016年2月4日, 截止2016年2月11日17时. 保留7题和全部子问. 原件由MIT课程组提供; 未取得官方解答. 以下均为\textbf{AI独立解答}. 第1、2题使用原题Gaussian单位, $e^2$已经包含SI中相应库仑常数; 不直接将SI电荷数值代入这些式子.

\begin{exercise}
经典氢原子中, 电子绕固定质子作近圆周运动. 原题给$V=-e^2/r$, 库仑力大小$e^2/r^2$, 及Larmor辐射功率$\mathrm dE/\mathrm dt=-2e^2a^2/(3c^3)$.
(a) 求每周损失能量$\Delta E$与动能$K$之比, 说明何时可把轨道视为逐渐缩小的圆.
(b) 在半径50 pm、1 pm、1 fm时比较$v$与$c$.
(c) 从50 pm缩到1 pm需多久? 非相对论近似是否合理? 形式上把终点改为1 fm有多大改变?
(d) $r\to0$时能量如何变化? 有无下界? （10分）
\end{exercise}
\begin{solution}
圆周平衡给$mv^2/r=e^2/r^2$, 所以$v^2=e^2/(mr)$、$K=e^2/(2r)$、$E=-e^2/(2r)$. 一周时间为$2\pi r/v$, 加速度为$v^2/r$. 辐射功率乘周期后,
$\Delta E=4\pi e^2v^3/(3c^3r)$, 故$\Delta E/K=(8\pi/3)(v/c)^3$. 当$v\ll c$, 每周能量损失很小, 近圆轨道慢变是自洽的局部近似.

令$r_e=e^2/(mc^2)=2.818\,\mathrm{fm}$, 则$v/c=\sqrt{r_e/r}$. 三个半径给约$0.00751$、$0.0531$和$1.68$. 最后一个超过1, 是非相对论模型失效的信号, 不能解释为电子实际超光速.

由$\dot E=(e^2/2r^2)\dot r$和辐射式, $\dot r=-4e^4/(3m^2c^3r^2)$. 积分得
\[
t=\frac{m^2c^3}{4e^4}(r_i^3-r_f^3)
=\frac{r_i^3-r_f^3}{4cr_e^2}.
\]
取50 pm和1 pm得$t\simeq1.31\times10^{-11}\,\mathrm s$. 到1 pm速度仍约$c$的5\%, 相对论修正小. 把终点形式上改为1 fm, 时间仅增加约原值的$8\times10^{-6}$, 因为积分受起始大半径主导. 但很小半径处的辐射与轨道模型已不自洽, 这个形式估值不能验证最后阶段的真实运动. 经典$E=-e^2/(2r)$无下界, $r\to0$趋$-\infty$, 正是稳定性失败.
\end{solution}

\begin{exercise}
电子绕固定重质子作圆周运动, $V=-e^2/r$.
(a) 给$K,V,E$间关系.
(b) 若角动量大小量子化为$L=n\hbar$, $n$为正整数, 求$E_n,r_n$, 用电子静能$E_e=mc^2$、约化Compton长度$\bar\lambda=\hbar/(mc)$、$\alpha=e^2/(\hbar c)$表达. （5分）
\end{exercise}
\begin{solution}
上题的圆周条件给$K=-V/2$, $E=V/2=-K$. 又$L^2=m^2v^2r^2=me^2r$, 故$r_n=n^2\hbar^2/(me^2)=n^2\bar\lambda/\alpha$, $E_n=-mc^2\alpha^2/(2n^2)=-E_e\alpha^2/(2n^2)$. 这是Bohr旧量子模型, 能级与非相对论氢谱同形, 但把真实电子态等同确定圆轨道不是现代量子力学的结论.
\end{solution}

\begin{exercise}
取$h=6.6\times10^{-34}\,\mathrm{J\,s}$, $k_B=1.38\times10^{-23}\,\mathrm{J/K}$.
(a) 求下列de Broglie波长并解释日常为何不易直接看到它: 2000 kg汽车以22 m/s运动; 10 g弹珠以10 cm/s运动; 直径100 nm、质量1 fg烟粒在300 K具有平均平动动能$3k_BT/2$; $^{87}$Rb原子冷却到100 $\mu$K且同样取$3k_BT/2$.
(b) 光子$E=h\nu$: (i) 求400--700 nm可见光的频率和能量范围; (ii) 求2.5 GHz、300 W微波炉, 633 nm、10 mW激光, 850 MHz、0.25 W手机每秒发出光子数; (iii) 200 ml水升温10 K需多少该微波光子, 水密度1 g/ml、比热4 J/(g K)? (iv) 同功率下无线电与X射线何者更适合经典波近似? （10分）
\end{exercise}
\begin{solution}
非相对论$p=mv$, 热动能条件给$p=\sqrt{3mk_BT}$; 1 fg为$10^{-18}$ kg. 四种波长依次约为$1.50\times10^{-38}$ m、$6.6\times10^{-31}$ m、$5.92\times10^{-15}$ m和$2.70\times10^{-8}$ m. 最后一项使用$^{87}$Rb质量约$87\times1.66\times10^{-27}$ kg. 可见烟粒的几何直径与其中心质点波长是不同概念. 宏观物体波长很小, 又难保持不同路径相干, 日常难以分辨其量子干涉.

取$c=3.0\times10^8$ m/s. 可见光频率约$4.29\times10^{14}$到$7.50\times10^{14}$ Hz, 能量$2.83\times10^{-19}$到$4.95\times10^{-19}$ J, 约1.77--3.09 eV. 低波长端对应高频高能. 光子率为$\dot N=P/(h\nu)$, 三装置分别约$1.82\times10^{26}$、$3.20\times10^{16}$、$4.46\times10^{23}\,\mathrm{s^{-1}}$. 激光频率须先用$c/\lambda$计算.

水的吸收能量为$200\times4\times10=8000$ J, 每个微波光子$1.65\times10^{-24}$ J, 所以需约$4.85\times10^{27}$个, 假设这些能量全部被水吸收. 同功率和相同观察时间下, 低频光每光子能量更低、光子数更多, 相对的离散光子涨落较小, 通常更适合经典波统计近似. 实际经典性还依赖制备和相干性, 不能仅由频率对所有量子光态下断言.
\end{solution}

\begin{exercise}
写$z=a+\ii b=r\e^{\ii\theta}$.
(a) 用Taylor级数证明Euler公式.
(b) 写笛卡尔与极坐标参数的转换.
(c) 证明乘$\ii$等于转$90^\circ$, 给$\ii z$的实部, 并证明乘$\e^{\ii\phi}$等于转$\phi$.
(d) 对$z^*=a-\ii b$: (i) 是否有数同时为实数和纯虚数? (ii) 求$(z^*)^*$及$z^*$的极形式; (iii) 用$z,z^*$表达实虚部; (iv) 求$zz^*$的多种表达.
(e) 推出二倍角、三倍角及和角公式. （15分）
\end{exercise}
\begin{solution}
绝对收敛指数级数中偶次项为$\sum_j(-1)^j\theta^{2j}/(2j)!=\cos\theta$, 奇次项为$\ii\sum_j(-1)^j\theta^{2j+1}/(2j+1)!=\ii\sin\theta$, 故Euler公式成立. $a=r\cos\theta$, $b=r\sin\theta$, $r=\sqrt{a^2+b^2}$, 当$r>0$时$\theta=\operatorname{atan2}(b,a)\pmod{2\pi}$, 即按象限取角; $r=0$时角度任意.

$\ii=\e^{\ii\pi/2}$, 所以$\ii z=r\e^{\ii(\theta+\pi/2)}=-b+\ii a$, 实部$-b$. 一般乘相位只将$\theta$加$\phi$. 同时实且纯虚的唯一数为0; $(z^*)^*=z$, $z^*=r\e^{-\ii\theta}$, $\operatorname{Re}z=(z+z^*)/2$, $\operatorname{Im}z=(z-z^*)/(2\ii)$, $zz^*=a^2+b^2=r^2=|z|^2$.

由$\e^{2\ii\theta}=(\cos\theta+\ii\sin\theta)^2$比较实虚部, 得$\cos2\theta=\cos^2\theta-\sin^2\theta$, $\sin2\theta=2\sin\theta\cos\theta$. 立方比较给$\cos3\theta=\cos^3\theta-3\cos\theta\sin^2\theta$, $\sin3\theta=3\sin\theta\cos^2\theta-\sin^3\theta$. 最后$\e^{\ii(A+B)}=\e^{\ii A}\e^{\ii B}$给$\cos(A+B)=\cos A\cos B-\sin A\sin B$, $\sin(A+B)=\sin A\cos B+\cos A\sin B$.
\end{solution}

\begin{exercise}
光子与自由电子碰撞, 解释光子为何不能被完全吸收. （5分）
\end{exercise}
\begin{solution}
在电子初始静止系, 设光子能量$\epsilon>0$. 若最终只有一个电子, 动量守恒给$p_f=\epsilon/c$, 能量守恒给$E_f=mc^2+\epsilon$. 但电子质量壳条件要求$E_f^2=p_f^2c^2+m^2c^4$, 代入得到$2mc^2\epsilon=0$, 矛盾. 因而至少还须有散射光子或其他物体带走动量; 束缚电子可以与原子核共同分担反冲, 不受这个自由二体约束.
\end{solution}

\begin{exercise}
Mach--Zehnder干涉仪入射态为$(\alpha,\beta)^{\mathsf T}$, $|\alpha|^2+|\beta|^2=1$, 探测器为$D_0,D_1$.
(a) 在下臂放相位$\delta_l$, 求$P_0,P_1$.
(b) 只在上臂放$\delta_u$.
(c) 两个相位器同时放入. （10分）
\end{exercise}
\begin{solution}
使用原讲义Lecture 2的两个分束器约定
$B_1=\frac1{\sqrt2}\begin{pmatrix}-1&1\\1&1\end{pmatrix}$,
$B_2=\frac1{\sqrt2}\begin{pmatrix}1&1\\1&-1\end{pmatrix}$.
上、下臂相位矩阵为$D=\operatorname{diag}(\e^{\ii\delta_u},\e^{\ii\delta_l})$. 输出为$B_2DB_1(\alpha,\beta)^{\mathsf T}$. 令$\delta=\delta_l-\delta_u$, 去掉不影响概率的整体$\e^{\ii\delta_u}$, 得
\begin{align*}
P_0(\delta)&=\frac14|\alpha(\e^{\ii\delta}-1)+\beta(1+\e^{\ii\delta})|^2,\\
P_1(\delta)&=\frac14|-\alpha(1+\e^{\ii\delta})+\beta(1-\e^{\ii\delta})|^2.
\end{align*}
(a)代$\delta=\delta_l$, (b)代$\delta=-\delta_u$, (c)代$\delta=\delta_l-\delta_u$. 因而两相位都相同不改变结果. $\delta=0$时$P_0=|\beta|^2$, $P_1=|\alpha|^2$, 与原讲义输出$(\beta,-\alpha)$一致. 不能把后一测试中两个相同分束器的约定直接套入本题.
\end{solution}

\begin{exercise}
Elitzur--Vaidman思想实验采用原讲义的平衡干涉仪, 从下方入射, 下臂放待检光子吸收器.
(a) 重复测试直至有足够把握判断状态, 有效吸收器有多少比例可在未触发吸收的情况下被认证?
(b) 80\%为失效吸收器. 随机取一个, 发入一个光子, 在$D_0$探测到. 该吸收器失效的条件概率是多少? 原题以“炸弹”命名这些理想化吸收器. （10分）
\end{exercise}
\begin{solution}
失效时两臂畅通,$D_0$概率1,$D_1$概率0. 有效时一次测试中, 吸收概率$1/2$, $D_0$与$D_1$各$1/4$. $D_1$事件认证有效且没有吸收; $D_0$不确定, 可重试. 对有效者, 认证比例是$\sum_{j\ge0}(1/4)^j(1/4)=1/3$. 这是重复次数趋大的理想极限; 有限次连续$D_0$不能逻辑上完全证明失效, 只能按先验降低有效概率.

Bayes规则给$\Pr(\text{失效}\mid D_0)=0.8\times1/[0.8\times1+0.2\times(1/4)]=16/17$. 将条件概率直接写成80\%会忽略此次测量提供的信息.
\end{solution}
